\documentclass{article}
\usepackage[T1]{fontenc}
\usepackage{lmodern}
\usepackage{graphicx}
\usepackage{amsmath,amssymb}
\usepackage[a4paper,margin=2cm]{geometry}
\usepackage[absolute,overlay]{textpos}
\setlength{\TPHorizModule}{1mm}
\setlength{\TPVertModule}{1mm}
\usepackage[indonesian]{babel}
\usepackage{hyperref}
\setlength{\emergencystretch}{2em}
% unit: Halaman 4

\begin{document}

% === Halaman 4 (llm) ===

\begin{itemize}
    \item Dalam teori algoritma, persoalan \textit{knapsack} termasuk ke dalam kelompok \textit{NP-complete}.
    \item Persoalan yang termasuk \textit{NP-complete} tidak dapat dipecahkan dalam orde waktu polinomial.
    \item Ide dasar dari algoritma kriptografi \textit{knapsack} adalah mengkodekan pesan sebagai rangkaian solusi dari persoalan \textit{knapsack}.
    \item Setiap bobot $w_i$ di dalam persoalan \textit{knapsack} merupakan kunci rahasia, sedangkan bit-bit plainteks menyatakan $b_i$.
\end{itemize}

\end{document}
